\section{The logarithmic consequence of a Bell overlap}

\begin{definition}[The overlap remainder]
    \label{def:bell_overlap_remainder}
    \lean{Real.bellOverlapRemainder}
    \leanok
    For natural numbers $d,K,k$ and $c\in\R$, put
    \begin{align*}
        R_{d,c,K}(k)
        =\frac{\log4+(d^2+cK)\log(k+2)}{k}.
    \end{align*}
    The estimates below concern positive $k$.
\end{definition}

\begin{theorem}[A logarithmic upper bound from the overlap]
    \label{thm:bell_overlap_log_bound}
    \lean{Real.neg_log_div_le_of_bell_overlap_bound}
    \leanok
    \uses{def:bell_overlap_remainder}
    Let $d,k$ be positive integers, $K$ a nonnegative integer, $z>0$,
    $c\ge0$, and $0\le2s\le1$. Let $t,B,r\in\R$. Suppose that
    \begin{align*}
        \frac{(z/d^2)^k}{4(k+1)^{d^2}}
        \le\bigl((k+2)^{cK}\exp(4tKkB)\bigr)^{2s}r.
    \end{align*}
    Then $r>0$ and
    \begin{align*}
        -\frac{\log r}{k}
        \le2\log d-\log z+8stKB+R_{d,c,K}(k).
    \end{align*}
    The remainder is independent of $s,t,B,z$.
\end{theorem}
\begin{proof}
    \leanok
    The left side of the assumed inequality and the coefficient of $r$
    are positive. Thus $r>0$. Taking logarithms, expanding the products
    and powers, and dividing by $k$ gives the stated leading terms with
    remainder
    \begin{align*}
        \frac{\log4+d^2\log(k+1)+2scK\log(k+2)}{k}.
    \end{align*}
    Use $\log(k+1)\le\log(k+2)$ and $2s\le1$ to bound this remainder
    by $R_{d,c,K}(k)$. This is the scalar step in
    \cite[\texttt{07-comparators.tex}, lines 153--163 and 332--381]{OpenAI2026AreaLaw}.
    In the physical comparison, $r$ is the squared inverse-power norm
    of the original auxiliary-label prevector; the assumed overlap
    inequality must be obtained for that same vector.
\end{proof}

\begin{theorem}[Vanishing of the overlap remainder]
    \label{thm:bell_overlap_remainder_tendsto}
    \lean{Real.tendsto_bellOverlapRemainder}
    \leanok
    \uses{def:bell_overlap_remainder}
    For every fixed $d,K\in\N$ and $c\in\R$,
    \begin{align*}
        \lim_{k\to\infty}R_{d,c,K}(k)=0.
    \end{align*}
\end{theorem}
\begin{proof}
    \leanok
    Both $1/k$ and $\log(k+2)/k$ tend to zero. The coefficients
    $\log4$ and $d^2+cK$ are fixed. This gives the vanishing error
    used in the replica limit in
    \cite[\texttt{07-comparators.tex}, lines 659--667]{OpenAI2026AreaLaw}.
\end{proof}
